\begin{resumo}
Órgãos federais acompanham políticas públicas com dados espalhados por
sistemas que foram construídos e são geridos separadamente, e poucos deles
compartilham uma chave que uma junção possa usar diretamente. Os
identificadores mudam de formato, precisam ser combinados ou aparecem apenas
em texto livre. No Gov Hub, plataforma de dados de código aberto desenvolvida
pela Universidade de Brasília com ministérios federais, apenas uma das 29
condições de junção entre fontes do código de produção é uma igualdade direta
entre colunas sem transformação. Os \textit{schema matchers} não produzem
essas transformações, e um modelo de linguagem de grande porte (LLM) que
escreve o SQL da junção pode devolver uma junção errada que executa sem erro.
Este trabalho propõe um sistema em que o LLM apenas escolhe a chave, a partir
de evidências coletadas antes e de um catálogo fechado de operações, ou se
abstém, enquanto código determinístico gera e verifica a junção em dbt. O
mesmo princípio, que chamamos de agência limitada, é aplicado ao agente de IA
de código que constrói o sistema por \textit{Spec-Driven Development}. O
método é a \textit{Design Science Research} com estudo de caso no repositório
do Ministério da Cultura no Gov Hub. A avaliação planejada compara o sistema
com \textit{matchers} clássicos, regras e outros usos do LLM num benchmark de
pares reais e perturbados, compara modelos abertos e proprietários, mede o
efeito sobre analistas, aplica o sistema a pares de fontes do ministério
ainda não integrados e compara o trabalho do agente com e sem especificação
por testes escondidos dele. Esta primeira etapa apresenta o problema, o
desenho do sistema, o protocolo do estudo e os resultados preliminares.

 \vspace{\onelineskip}

 \noindent
 \textbf{Palavras-chave}: integração de dados. modelos de linguagem.
 \textit{spec-driven development}. dados abertos governamentais.
 \textit{schema matching}.
\end{resumo}
